%%%%%%%%%%%%%%%%%%%%%%%%%%%%%%%%%%%%%%%%%%%%%%%%%%%%%%%%%%%%%%%%%%%%%%%%%%%%%%%%
% SECURITY ANALYSIS
%%%%%%%%%%%%%%%%%%%%%%%%%%%%%%%%%%%%%%%%%%%%%%%%%%%%%%%%%%%%%%%%%%%%%%%%%%%%%%%%

\chapter[Security Analysis]{\IconShield\ Security Analysis}

This chapter presents an overall assessment of the security posture of the \textbf{\ProjectName} based on the generated Software Bill of Materials (SBOM) and the Software Composition Analysis (SCA) performed using Grype and the ZeroWatch Audit Scan. While the SBOM generated by Syft provides a complete inventory of software components, Grype and the ZeroWatch Audit Scan evaluate those components against multiple vulnerability intelligence sources to identify publicly disclosed security issues and available remediation guidance.

Rather than focusing on individual vulnerabilities, this chapter evaluates the overall security readiness of the software supply chain, the effectiveness of the adopted Software Composition Analysis process, and the practices required to maintain a secure dependency ecosystem. Detailed vulnerability information is presented separately in the following chapter.

\vspace{0.4cm}

%%%%%%%%%%%%%%%%%%%%%%%%%%%%%%%%%%%%%%%%%%%%%%%%%%%%%%%%%%%%%%%%%%%%%%%%%%%%%%%%
% SECURITY OVERVIEW
%%%%%%%%%%%%%%%%%%%%%%%%%%%%%%%%%%%%%%%%%%%%%%%%%%%%%%%%%%%%%%%%%%%%%%%%%%%%%%%%

\section{Security Overview}

\begin{table}[H]
\centering
\small
\renewcommand{\arraystretch}{1.3}

\begin{tabular}{p{5cm}p{8cm}}
\toprule
\rowcolor{ZWPrimary}
\textbf{\color{white}Security Metric} & \textbf{\color{white}Status} \\
\midrule
SBOM Generated & Yes \\
SBOM Standard & \SbomFormat\ \SbomSpecVersion \\
SBOM Generator & \SbomGenerator \\
Vulnerability Assessment & Completed using \VulScanner \\
Component Inventory & Available \\
Dependency Relationships & Available \\
Package URLs (PURLs) & Available \\
License Metadata & Available \\
Software Composition Analysis & Completed \\
Available Security Fixes & Yes \\
\bottomrule
\end{tabular}
\caption{Security Analysis Overview}
\end{table}

\vspace{0.4cm}

%%%%%%%%%%%%%%%%%%%%%%%%%%%%%%%%%%%%%%%%%%%%%%%%%%%%%%%%%%%%%%%%%%%%%%%%%%%%%%%%
% SECURITY OBJECTIVES
%%%%%%%%%%%%%%%%%%%%%%%%%%%%%%%%%%%%%%%%%%%%%%%%%%%%%%%%%%%%%%%%%%%%%%%%%%%%%%%%

\section{Security Objectives}

The combined SBOM generation and Software Composition Analysis support the following security objectives:
\begin{itemize}
\item Establish complete visibility into software components.
\item Improve software supply chain transparency.
\item Identify vulnerable third-party dependencies.
\item Support risk-based vulnerability prioritization.
\item Enable continuous dependency monitoring.
\item Improve software governance.
\item Strengthen software release assurance.
\item Support regulatory and customer compliance requirements.
\end{itemize}

\vspace{0.4cm}

%%%%%%%%%%%%%%%%%%%%%%%%%%%%%%%%%%%%%%%%%%%%%%%%%%%%%%%%%%%%%%%%%%%%%%%%%%%%%%%%
% SOFTWARE SUPPLY CHAIN SECURITY
%%%%%%%%%%%%%%%%%%%%%%%%%%%%%%%%%%%%%%%%%%%%%%%%%%%%%%%%%%%%%%%%%%%%%%%%%%%%%%%%

\section{Software Supply Chain Security}

Modern software applications depend heavily on third-party open-source libraries, making software supply chain security an essential aspect of secure software development. Every dependency introduces potential security, operational, licensing, and maintenance risks that must be continuously evaluated throughout the software lifecycle.

The Software Bill of Materials generated for the \ProjectName\ provides complete visibility into software dependencies, while Grype and the ZeroWatch Audit Scan evaluate these components against multiple vulnerability databases, including GitHub Security Advisories and the National Vulnerability Database (NVD). Together, these technologies enable rapid identification of vulnerable components, improve traceability, and support timely remediation.

\vspace{0.4cm}

%%%%%%%%%%%%%%%%%%%%%%%%%%%%%%%%%%%%%%%%%%%%%%%%%%%%%%%%%%%%%%%%%%%%%%%%%%%%%%%%
% SOFTWARE COMPOSITION ANALYSIS
%%%%%%%%%%%%%%%%%%%%%%%%%%%%%%%%%%%%%%%%%%%%%%%%%%%%%%%%%%%%%%%%%%%%%%%%%%%%%%%%

\section{Software Composition Analysis}

Software Composition Analysis (SCA) combines software inventory information with vulnerability intelligence to evaluate the security of third-party dependencies.

For the \ProjectName, the assessment follows the workflow below:
\begin{enumerate}
\item Generate a CycloneDX Software Bill of Materials using Syft.
\item Identify all direct and transitive software dependencies.
\item Correlate software components with public vulnerability databases using Grype and the ZeroWatch Audit Scan.
\item Review identified security advisories.
\item Prioritize remediation based on vulnerability severity.
\item Reassess dependencies after package upgrades.
\end{enumerate}

This workflow enables continuous monitoring of the software supply chain while maintaining an accurate inventory of software components.

\vspace{0.4cm}

%%%%%%%%%%%%%%%%%%%%%%%%%%%%%%%%%%%%%%%%%%%%%%%%%%%%%%%%%%%%%%%%%%%%%%%%%%%%%%%%
% POTENTIAL SECURITY RISK AREAS
%%%%%%%%%%%%%%%%%%%%%%%%%%%%%%%%%%%%%%%%%%%%%%%%%%%%%%%%%%%%%%%%%%%%%%%%%%%%%%%%

\section{Potential Security Risk Areas}

\begin{table}[H]
\centering
\small
\renewcommand{\arraystretch}{1.3}

\begin{tabular}{p{6cm}p{5cm}}
\toprule
\rowcolor{ZWPrimary}
\textbf{\color{white}Risk Area} & \textbf{\color{white}Potential Impact} \\
\midrule
Outdated Dependencies & High \\
Known Vulnerabilities & High \\
Dependency Confusion & High \\
Package Tampering & High \\
Unknown License Metadata & Medium \\
Abandoned Packages & Medium \\
Transitive Dependencies & Medium \\
Supply Chain Attacks & High \\
\bottomrule
\end{tabular}
\caption{Potential Software Supply Chain Risks}
\end{table}

\vspace{0.4cm}

%%%%%%%%%%%%%%%%%%%%%%%%%%%%%%%%%%%%%%%%%%%%%%%%%%%%%%%%%%%%%%%%%%%%%%%%%%%%%%%%
% CURRENT SECURITY POSTURE
%%%%%%%%%%%%%%%%%%%%%%%%%%%%%%%%%%%%%%%%%%%%%%%%%%%%%%%%%%%%%%%%%%%%%%%%%%%%%%%%

\section{Current Security Posture}

The \ProjectName\ demonstrates a mature software supply chain security approach by combining Software Bill of Materials generation with continuous Software Composition Analysis.

The generated SBOM provides complete visibility into third-party software components, while the vulnerability assessment verifies their current security status against publicly disclosed vulnerability databases. This approach enables proactive identification of security issues before software deployment and supports informed risk management decisions.

Maintaining this security posture requires periodic dependency updates, continuous vulnerability monitoring, regeneration of the SBOM after dependency changes, and regular reassessment of third-party software throughout the software development lifecycle.

\vspace{0.4cm}

%%%%%%%%%%%%%%%%%%%%%%%%%%%%%%%%%%%%%%%%%%%%%%%%%%%%%%%%%%%%%%%%%%%%%%%%%%%%%%%%
% RECOMMENDED SECURITY PRACTICES
%%%%%%%%%%%%%%%%%%%%%%%%%%%%%%%%%%%%%%%%%%%%%%%%%%%%%%%%%%%%%%%%%%%%%%%%%%%%%%%%

\section{Recommended Security Practices}

To maintain a secure software supply chain, the following practices are recommended:
\begin{itemize}
\item Generate a new SBOM for every software release.
\item Perform Grype and ZeroWatch Audit vulnerability assessments during every release cycle.
\item Upgrade dependencies with available security fixes.
\item Continuously monitor newly published security advisories.
\item Maintain integrated SBOM generation and vulnerability scanning in the CI/CD pipeline.
\item Verify package authenticity before deployment.
\item Maintain software provenance and dependency traceability.
\item Periodically review dependency health.
\end{itemize}

\vspace{0.4cm}

%%%%%%%%%%%%%%%%%%%%%%%%%%%%%%%%%%%%%%%%%%%%%%%%%%%%%%%%%%%%%%%%%%%%%%%%%%%%%%%%
% SECURITY BENEFITS
%%%%%%%%%%%%%%%%%%%%%%%%%%%%%%%%%%%%%%%%%%%%%%%%%%%%%%%%%%%%%%%%%%%%%%%%%%%%%%%%

\section{Security Benefits of SBOM and Vulnerability Assessment}

The combination of the generated Software Bill of Materials (SBOM), Grype, and ZeroWatch Audit Software Composition Analysis provides the following security benefits:
\begin{itemize}
\item Complete software component inventory.
\item Improved software traceability.
\item Automated identification of vulnerable software components.
\item Continuous Software Composition Analysis.
\item Faster vulnerability remediation.
\item Better software release readiness.
\item Improved third-party software governance.
\item Enhanced software supply chain transparency.
\item Simplified compliance reporting.
\item Improved incident response capabilities.
\end{itemize}

\vspace{0.4cm}

%%%%%%%%%%%%%%%%%%%%%%%%%%%%%%%%%%%%%%%%%%%%%%%%%%%%%%%%%%%%%%%%%%%%%%%%%%%%%%%%
% SUMMARY
%%%%%%%%%%%%%%%%%%%%%%%%%%%%%%%%%%%%%%%%%%%%%%%%%%%%%%%%%%%%%%%%%%%%%%%%%%%%%%%%

\section{Summary}

The Software Bill of Materials generated using Syft, together with the Grype and ZeroWatch Audit Software Composition Analysis, provides a comprehensive framework for managing software supply chain security within the \ProjectName. The SBOM delivers complete visibility into software dependencies, while vulnerability assessments enable continuous evaluation of the project's security posture.

Together, these capabilities improve software governance, support regulatory compliance, facilitate vulnerability management, and strengthen the overall security of the application throughout its development lifecycle. The detailed vulnerability findings, affected packages, security advisories, and remediation recommendations are presented in the following chapter.